\documentclass[12pt,a4paper]{article}
\usepackage[utf8]{inputenc}
\usepackage[spanish]{babel}
\usepackage{graphicx}
\usepackage{hyperref}
\usepackage{geometry}
\usepackage{booktabs}
\usepackage{listings}
\usepackage{xcolor}
\usepackage{forest}

\geometry{margin=2.5cm}

\lstset{
    basicstyle=\small\ttfamily,
    breaklines=true,
    frame=single,
    backgroundcolor=\color{gray!10}
}

\title{Fase 2: Implementacion ETL y Sitio Web\\Sistema de Ruteo Resiliente}
\author{Felipe Aros}
\date{Noviembre 2025}

\begin{document}

\maketitle

\section{Estructura del Proyecto}

El proyecto sigue una arquitectura modular con las siguientes carpetas:

\begin{verbatim}
ruteo/
├── infraestructura/     # ETL red vial OSM
├── metadata/            # ETL datos ambientales
├── amenazas/            # ETL fuentes de amenazas
├── sql/                 # Esquema y funciones BD
├── web/                 # Frontend Next.js
├── docker/              # Contenedores
├── main.py              # Orquestador ETL
└── load_to_supabase.py  # Carga a BD
\end{verbatim}

\section{Diagrama de Base de Datos}

El esquema utiliza PostgreSQL con las extensiones PostGIS y pgRouting:

\begin{lstlisting}[language=SQL]
-- Extensiones requeridas
CREATE EXTENSION IF NOT EXISTS postgis;
CREATE EXTENSION IF NOT EXISTS pgrouting;

-- Tabla de nodos (42,534 registros)
CREATE TABLE infra_nodos (
    id BIGSERIAL PRIMARY KEY,
    geom GEOMETRY(Point, 4326) NOT NULL,
    osm_id BIGINT,
    tags JSONB,
    p_fallo_nodo DOUBLE PRECISION DEFAULT 0.0
);

-- Tabla de aristas (89,021 registros)
CREATE TABLE infra_aristas (
    id BIGSERIAL PRIMARY KEY,
    source BIGINT REFERENCES infra_nodos(id),
    target BIGINT REFERENCES infra_nodos(id),
    geom GEOMETRY(LineString, 4326) NOT NULL,
    length_m DOUBLE PRECISION,
    cost DOUBLE PRECISION,
    p_fallo_arista DOUBLE PRECISION DEFAULT 0.0,
    highway TEXT
);

-- Indices espaciales GiST
CREATE INDEX ON infra_nodos USING GIST(geom);
CREATE INDEX ON infra_aristas USING GIST(geom);
\end{lstlisting}

\section{Carpeta Infraestructura}

\subsection{Archivo: osm\_infra\_etl.py}

Automatiza la extraccion de la red vial desde OpenStreetMap:

\begin{lstlisting}[language=Python]
# Consulta Overpass API
OVERPASS_QUERY = """
[out:json][timeout:300];
area["name"="Santiago"]["admin_level"="8"]->.a;
(
  way(area.a)["highway"];
  way(area.a)["highway"~"primary|secondary|tertiary"];
);
out geom;
"""
\end{lstlisting}

\textbf{Salida}: \texttt{infraestructura.geojson}

\subsection{Estructura del JSON}

\begin{lstlisting}[language=json]
{
  "type": "FeatureCollection",
  "features": [
    {
      "type": "Feature",
      "geometry": {
        "type": "LineString",
        "coordinates": [[-70.65, -33.45], ...]
      },
      "properties": {
        "osm_id": 123456,
        "highway": "primary",
        "name": "Av. Providencia"
      }
    }
  ]
}
\end{lstlisting}

\section{Carpeta Metadata}

\subsection{Archivo 1: elevacion\_etl.py}

Descarga datos de elevacion SRTM de NASA Earthdata.

\textbf{Salida}: \texttt{elevacion.geojson} con campos:
\begin{itemize}
    \item \texttt{elev\_m}: Elevacion en metros sobre el nivel del mar
    \item \texttt{source}: "SRTM"
\end{itemize}

\subsection{Archivo 2: lluvia\_etl.py}

Obtiene datos de precipitacion historica.

\textbf{Salida}: \texttt{lluvia.geojson} con campos:
\begin{itemize}
    \item \texttt{precip\_mm}: Precipitacion en milimetros
    \item \texttt{observed\_at}: Fecha de observacion
\end{itemize}

\subsection{Archivo 3: landcover\_etl.py}

Clasifica la cobertura del suelo para estimar permeabilidad.

\textbf{Salida}: \texttt{landcover.geojson} con campos:
\begin{itemize}
    \item \texttt{class}: Tipo de cobertura (urbano, vegetacion, agua)
    \item \texttt{confidence}: Nivel de confianza
\end{itemize}

\section{Carpeta Amenazas}

\subsection{Archivo 1: dga\_etl.py}

Extrae estaciones de monitoreo de la Direccion General de Aguas.

\textbf{Fuente}: Portal MMA (\url{https://lineasdebasepublicas.mma.gob.cl})

\textbf{Salida}: \texttt{dga\_estaciones.json}

\subsection{Archivo 2: inundaciones\_hist\_etl.py}

Procesa zonas de inundacion historica basadas en SERNAGEOMIN.

\textbf{Salida}: \texttt{inundaciones\_hist.json} con campos:
\begin{itemize}
    \item \texttt{geometry}: Poligono de zona afectada
    \item \texttt{source}: Fuente (ej: "Rio Mapocho")
    \item \texttt{severity}: Nivel de severidad
\end{itemize}

\subsection{Archivo 3: reportes\_ciudadanos\_etl.py}

Geocodifica reportes ciudadanos de inundaciones.

\textbf{Salida}: \texttt{reportes\_ciudadanos.json}

\section{Sitio Web con Leaflet}

El frontend se implementa con Next.js 14 y Leaflet:

\begin{lstlisting}[language=JavaScript]
// web/src/app/page.tsx
import { MapContainer, TileLayer, GeoJSON } from 'react-leaflet';

export default function HomePage() {
  return (
    <MapContainer center={[-33.45, -70.65]} zoom={12}>
      <TileLayer url="https://{s}.tile.openstreetmap.org/{z}/{x}/{y}.png" />
      <GeoJSON data={infraestructura} />
      <GeoJSON data={inundaciones} style={{color: 'red'}} />
      <GeoJSON data={ruta} style={{color: 'blue', weight: 5}} />
    </MapContainer>
  );
}
\end{lstlisting}

\section{Ruta con pgr\_dijkstra}

Ejemplo de ruta basica usando solo distancia como costo:

\begin{lstlisting}[language=SQL]
SELECT seq, edge, cost, agg_cost, ST_AsGeoJSON(geom)
FROM pgr_dijkstra(
    'SELECT id, source, target, length_m AS cost
     FROM infra_aristas',
    (SELECT id FROM infra_nodos ORDER BY geom <->
     ST_SetSRID(ST_MakePoint(-70.60, -33.42), 4326) LIMIT 1),
    (SELECT id FROM infra_nodos ORDER BY geom <->
     ST_SetSRID(ST_MakePoint(-70.65, -33.45), 4326) LIMIT 1)
) AS di
JOIN infra_aristas a ON di.edge = a.id;
\end{lstlisting}

\section{Archivo main.py}

Orquesta todo el proceso ETL:

\begin{lstlisting}[language=Python]
#!/usr/bin/env python3
"""Orquestador ETL principal"""

from infraestructura import osm_infra_etl
from metadata import elevacion_etl, lluvia_etl, landcover_etl
from amenazas import dga_etl, inundaciones_hist_etl, reportes_ciudadanos_etl

def main():
    print("1. Extrayendo infraestructura OSM...")
    osm_infra_etl.run()

    print("2. Procesando metadata...")
    elevacion_etl.run()
    lluvia_etl.run()
    landcover_etl.run()

    print("3. Procesando amenazas...")
    dga_etl.run()
    inundaciones_hist_etl.run()
    reportes_ciudadanos_etl.run()

    print("ETL completado!")

if __name__ == "__main__":
    main()
\end{lstlisting}

\section{Docker}

\begin{lstlisting}[language=yaml]
# docker/docker-compose.yml
version: '3.8'
services:
  app:
    build:
      context: ..
      dockerfile: docker/Dockerfile.app
    ports:
      - "3000:3000"
    environment:
      - SUPABASE_DB_HOST=${SUPABASE_DB_HOST}
      - SUPABASE_DB_PORT=6543
    volumes:
      - ..:/app
\end{lstlisting}

\end{document}
